\section{Source lock and exact-ledger rigidity}
\label{sec:rigidity}

Let $\At=\{F_{p_n}:n\ge1\}$ be the tensor atoms of the full-shift monoid,
listed by increasing topological entropy.  Unique factorization and strict
monotonicity of $\log n$ give the same order as
$2=p_1<p_2<p_3<\cdots$.  This is an intrinsic enumeration; no external
prime table is part of the definition.

The phase space of SD-C09 is the countable directed graph with vertex set
$\At$.  It has a loop at every vertex and two parallel edges from $p_n$ to
$p_{n+1}$.  The loop has roof $\log p_n$.  The two successor edges carry,
respectively, the source roof $\log p_n$ and target roof $\log p_{n+1}$,
each with scalar potential $-\log2$.  There are no decreasing edges.  The
function space is $\ell^2(\At)$; its transfer is defined in
\cref{sec:coupling}.  This freezes the object, grammar, roof, potential,
function space, and determinant convention before any spectral comparison.

\subsection{A universal one-sided gauge theorem}

We first correct and strengthen the previous stage's no-motion theorem.  Let
$G$ be any positive injective diagonal mass operator and let $K$ be bounded.
No commutation is assumed.

\begin{theorem}[Universal one-sided phase gauge]
\label{thm:one-sided-gauge}
For
\[
 A_t=G^{1/2+it}K,
 \qquad
 B_t=\begin{pmatrix}0&A_t\\A_t^*&0\end{pmatrix},
\]
one has
\[
 B_t=U_tB_0U_t^*,
 \qquad U_t=\operatorname{diag}(G^{it},I).
\]
Thus the spectrum, singular values, and every unitarily invariant regularized
determinant are independent of $t$, even when $[G,K]\ne0$.
\end{theorem}

\begin{proof}
Functional calculus gives $A_t=G^{it}A_0$.  Substitution into the two chiral
blocks yields the asserted conjugacy.
\end{proof}

Mass noncommutation is therefore necessary only after the phase has been
placed at both endpoints in a way that cannot be factored unitarily from one
side.

\subsection{Exact Euler ledgers force directed acyclicity}

The next result locates every scalar one-step escape.  For independent
commuting variables $x=(x_1,\ldots,x_N)$ put
$D(x)=\operatorname{diag}(x_1,\ldots,x_N)$.

\begin{theorem}[Exact-ledger DAG rigidity]
\label{thm:dag}
Let $K\in M_N(\CC)$ and assume
\[
 \det(I-zD(x)K)=\prod_{j=1}^N(1-zx_j)             \tag{3.1}
\]
as a polynomial identity in $z,x_1,\ldots,x_N$.  Then $K=I+N_0$, and the
directed support of $N_0$ contains no directed cycle.  Conversely, after a
permutation of the vertices, every strictly triangular $N_0$ satisfies
\cref{eq:dag-converse}:
\begin{equation}
 \det(I-zD(x)(I+N_0))=\prod_j(1-zx_j).
 \label{eq:dag-converse}
\end{equation}
The conclusion also follows from equality of all Dirichlet coefficients
after $x_j=p_j^{-s}$ for distinct tensor atoms.
\end{theorem}

\begin{proof}
Expanding the determinant by principal minors and comparing the coefficient
of $z^{|S|}\prod_{j\in S}x_j$ gives $\det K_{S,S}=1$ for every $S$.
In particular $K_{jj}=1$.  Write $K=I+N_0$ with zero diagonal.  The identity
\[
 \det(I+(N_0)_{S,S})=
 \sum_{T\subseteq S}\det (N_0)_{T,T}
\]
and induction on $|S|$ imply $\det(N_0)_{S,S}=0$ for every nonempty $S$.

If the support of $N_0$ contained a directed cycle, choose one of shortest
length with vertex set $C$.  It has no directed chord: a chord followed by
one of the two directed portions of the cycle would produce a shorter cycle.
The determinant of $(N_0)_{C,C}$ then contains exactly one nonzero
permutation term, the product around the cycle, contradicting its vanishing.
The support is a DAG and hence admits a topological ordering making $N_0$
strictly triangular.  The converse follows from triangularity.  Finally,
unique factorization makes distinct prime multisets linearly independent as
Dirichlet monomials, so the same coefficient comparison applies.
\end{proof}

The theorem covers arbitrary scalar signs and complex phases: exactness
does not leave a hidden recurrent signed solution.  A finite-dimensional
unitary cocycle does not evade it at the primitive-factor level, because a
mixed cycle labelled by an invertible $U$ contributes
$\det(I-wU)$, which is nonconstant.

\subsection{All-order graded cancellation is divisor-invisible}

\begin{theorem}[Trace-invisible cancellation]
\label{thm:trace-invisible}
Let $T=T^+\oplus T^-$ be a trace-class graded operator.  If
$\Str T^r=0$ for every $r\ge1$, then
\[
 \Ber(I-zT)=1
\]
throughout its Fredholm domain.  More generally, if $T^\pm\in S_q$ and
$\Tr(T^+)^r=\Tr(T^-)^r$ for every $r\ge q$, their relative
$q$-regularized determinant is one.
\end{theorem}

\begin{proof}
Use the convergent logarithmic expansions
\[
 \log\Ber(I-zT)=-\sum_{r\ge1}\frac{z^r}{r}\Str T^r,
\]
and, after the first $q-1$ subtractions, the analogous expansion beginning at
$r=q$.  Analytic continuation on the connected determinant domain preserves
the identity.
\end{proof}

Thus a chain-level sector may move spectrally while exactness cancels every
mixed trace, but that motion cannot generate zeros of the determinant whose
logarithm uses those same traces.  This is the precise ceiling that motivates
the acyclic rather than homological construction below.
